% TOP_PROBLEM: {"problemId": "problem.non-perturbative-construction-and-non-triviality-of-4d-euclidean-phi4-quantum-field-theory", "canonicalTitle": "Non-Perturbative Construction and Non-Triviality of 4D Euclidean Phi4 Quantum Field Theory", "releaseRank": 51, "primaryDomain": "mathematical_physics", "primaryDomainLabel": "Mathematical physics", "displayStatus": "reviewed_hold", "releaseStatus": "open"}
% !TeX program = lualatex
% Definitions and sources retained from research_batch_51_54.tex; research is in GPT_6_Astra_Ultra.
\documentclass[11pt]{article}
\usepackage[margin=25mm]{geometry}
\usepackage{fontspec}
\setmainfont{Latin Modern Roman}
\usepackage{amsmath,amssymb,mathtools}
\usepackage{unicode-math}
\setmathfont{Latin Modern Math}
\usepackage{enumitem,needspace}
\usepackage{xurl,hyperref}
\hypersetup{colorlinks=true,urlcolor=blue}
\setcounter{secnumdepth}{1}
\setlength{\parindent}{0pt}
\setlength{\parskip}{5pt}
\setlength{\emergencystretch}{3em}
\setlist[itemize]{leftmargin=3em,itemsep=5pt,topsep=4pt,parsep=0pt}
\allowdisplaybreaks
\newcommand{\N}{\mathbb{N}}
\newcommand{\Z}{\mathbb{Z}}
\newcommand{\Q}{\mathbb{Q}}
\newcommand{\R}{\mathbb{R}}
\newcommand{\C}{\mathbb{C}}
\newcommand{\F}{\mathbb{F}}
\newcommand{\A}{\mathbb{A}}
\newcommand{\PP}{\mathbb P}
\newcommand{\E}{\mathbb{E}}
\newcommand{\eps}{\varepsilon}
\newcommand{\dd}{\,\mathrm d}
\newcommand{\sref}[2]{\hyperlink{source:#1:#2}{[S#2]}}
\newcommand{\eref}[2]{\hyperlink{extra:#1:#2}{[E#2]}}
% Turn the next switch off for a shorter reading copy. The quotations preserve
% every exactTarget field, including qualifications omitted from a short summary.
\newif\ifshowcatalogue
\showcataloguetrue
\newcommand{\cataloguescope}[1]{\ifshowcatalogue
 \paragraph{Exact scope in the supplied catalogue.}
 \begin{quote}\small #1\end{quote}\fi}
\begin{document}
\setcounter{section}{50}
% ================================================================
% problemId: problem.non-perturbative-construction-and-non-triviality-of-4d-euclidean-phi4-quantum-field-theory
% source releaseRank: 51; source releaseStatus: open
% exactTarget SHA-256: f31f95cf191e6998f9d82a2d67885af59277b6579c9e6ee2e8b6d57c561713eb
\clearpage
\section{\texorpdfstring{Non-Perturbative Construction and Non-Triviality of 4D Euclidean Phi4 Quantum Field Theory}{Non-Perturbative Construction and Non-Triviality of 4D Euclidean Phi4 Quantum Field Theory}}\label{51}
\subsection{Definitions and mathematical statement}
A Euclidean scalar field is a probability law on generalized functions on $\R^4$, or equivalently a compatible collection of Schwinger correlation distributions. The formal quartic action is
\[
 S(\phi)=\int_{\R^4}\left(\tfrac12|\nabla\phi|^2+\tfrac12m^2\phi^2+\lambda\phi^4\right)dx.
\]
This expression by itself does not define the continuum measure. A construction must specify regularization, renormalization, and a limit satisfying the Osterwalder--Schrader conditions, including Euclidean invariance, reflection positivity, symmetry, and regularity. Non-Gaussian means that the correlations are not all determined by the two-point function through Wick's rule. \emph{Source-fit warning:} the cited Aizenman--Duminil-Copin theorem proves Gaussian limits for its four-dimensional ferromagnetic lattice Ising/$\phi^4$ class. The requested non-Gaussian construction cannot use that class of limits; the catalogue does not specify an alternative class. The record is retained, not certified as a well-posed currently open construction problem.
\par\smallskip\textit{Formulation and concept references:} \sref{51}{1}, \eref{51}{1}.
\cataloguescope{Construct a non-trivial, non-Gaussian Euclidean scalar quantum field theory with quartic interaction (phi\textasciicircum{}4)\_4 in four spacetime dimensions satisfying the Osterwalder-Schrader axioms.}
\subsection{Short English statement}
Can the proposed four-dimensional quartic scalar theory be constructed as a genuinely non-Gaussian continuum field, outside the class already ruled out by the cited triviality theorem?
\subsection{Sources}
\begin{itemize}\raggedright
\item[\textnormal{[S1]}]\hypertarget{source:51:1}{} M. Aizenman, H. Duminil-Copin, Ann. of Math. 194(1):163-235, 2021.
\url{https://doi.org/10.4007/annals.2021.194.1.3}.
{\small\textit{Catalogue formulation source.}}
\item[\textnormal{[S2]}]\hypertarget{source:51:2}{} Marginal triviality of the scaling limits of critical 4D Ising and phi4 models.
\url{https://annals.math.princeton.edu/2021/194-1/p03}.
{\small\textit{Catalogue research/status reference; not a proof endorsement.}}
\item[\textnormal{[S3]}]\hypertarget{source:51:3}{} The non triviality of a Phi4 model, III, the Osterwalder-Schrader Positivity.
\url{https://arxiv.org/abs/2101.11980}.
{\small\textit{Catalogue research/status reference; not a proof endorsement.}}
\item[\textnormal{[S4]}]\hypertarget{source:51:4}{} Triviality in a Non-Perturbative Second-Order Mean-Field Theory for phi4.
\url{https://arxiv.org/abs/2608.05998}.
{\small\textit{Catalogue research/status reference; not a proof endorsement.}}
\item[\textnormal{[E1]}]\hypertarget{extra:51:1}{} Michael Aizenman and Hugo Duminil-Copin, Marginal triviality of the scaling limits of critical 4D Ising and phi4 models, Annals of Mathematics 194 (2021), 163--235; arXiv:1912.07973.
\url{https://arxiv.org/abs/1912.07973}.
{\small\textit{Added primary source: Source-fit warning and proved triviality for the paper's model class. Consulted 24 September 2026.}}
\item[\textnormal{[E2]}]\hypertarget{extra:51:2}{} Romain Panis, \emph{Triviality of the scaling limits of critical Ising and $\phi^4$ models with effective dimension at least four}, arXiv:2309.05797; Annals of Probability, DOI 10.1214/25-AOP1782. Assumptions (A1)--(A6), Theorems 1.12--1.13 and Corollary 1.14.
\url{https://arxiv.org/abs/2309.05797}.
{\small\textit{Added research source: Quantitative triviality for the specified long-range model class; not a no-go theorem for arbitrary continuum constructions. Consulted 24 September 2026.}}
\item[\textnormal{[E3]}]\hypertarget{extra:51:3}{} Francesca Arici, Daniel Becker, Chris Ripken, Frank Saueressig and Walter D.\ van Suijlekom, \emph{Reflection positivity in higher derivative scalar theories}, arXiv:1712.04308 (2017), Theorem 3.7; DOI 10.1063/1.5027231.
\url{https://arxiv.org/abs/1712.04308}.
{\small\textit{Added research source: Rational-covariance reflection-positivity criterion. The negative reflection form used here is also derived explicitly in the attempt. Consulted 24 September 2026.}}
\item[\textnormal{[E4]}]\hypertarget{extra:51:4}{} Majdouline Borji, \emph{Triviality in a Non-Perturbative Second-Order Mean-Field Theory for $\phi^4$}, arXiv:2608.05998 (August 2026), concluding discussion of the ultraviolet remainder.
\url{https://arxiv.org/abs/2608.05998}.
{\small\textit{Added research source: Mean-field result and its stated limitation; no assertion of full-theory triviality is imported. Consulted 24 September 2026.}}
\item[\textnormal{[E5]}]\hypertarget{extra:51:5}{} Marietta Manolessou, \emph{The ``non triviality'' of a $\Phi^4_4$ model, III, the ``Osterwalder--Schrader Positivity''}, arXiv:2101.11980 (2021), Theorem 3.2.
\url{https://arxiv.org/abs/2101.11980}.
{\small\textit{Added research source: Claimed positivity result. The preceding hierarchy construction and the full construction claim were not independently verified in this notebook; not a proof endorsement. Consulted 24 September 2026.}}
% END RESEARCH ADDITION REFERENCES-51
\end{itemize}

\end{document}
